\begin{textsample}{POS Dim 3 – rs – Score 15.00 – rs/rs\_inf\_14.txt}  \label{ex:f3_pos_228}
11.3 CAMPO DE EXPERIÊNCIAS : TRAÇOS, \textbf{SONS}, CORES E FORMAS ( TS ) O Referencial Curricular Gaúcho, baseado no conceito deste Campo proposto pela BNCC (2017)[^7 ], destaca experiências nas quais as \textbf{crianças} tenham oportunidade de perceber o ambiente como um composto de traços, \textbf{sons}, cores e formas, oferecendo condições para sentirem a consistência da terra ou \textbf{areia}, criar misturas, colecionar coisas, modelar com argila, criar \textbf{tintas}, explorar formas coloridas, \textbf{texturas}, sabores, \textbf{sons} e também silêncios e um espaço acolhedor \textbf{cheio} de visualidades e sonoridades, promovendo o desenvolvimento da expressividade e da criatividade \textbf{infantil} e abrindo caminhos para o desenvolvimento de sua afetividade.
O Campo de Experiências Traços, \textbf{Sons}, Cores e Formas refere -se às múltiplas linguagens artísticas que envolvem expressão : música, dança, escultura, cinema e teatro.
Cada linguagem é constituída por diferentes elementos – imagens, cores, \textbf{sons} e traços – e utilizadas pelas \textbf{crianças} para se comunicar, expressar e interagir com o meio.
Cada linguagem possibilita a ampliação de vivências que consideram a expressão criativa das \textbf{crianças}, a partir das condições propostas e oferecidas pelo professor.
Um trabalho que exige planejamento, \textbf{acompanhamento}, avaliação e, muitas vezes, intervenções.
Nessa perspectiva, o trabalho com as linguagens artísticas não visa a formação de artistas, mas, auxiliar, através das diferentes linguagens e da arte, na formação de \textbf{crianças} sensíveis ao mundo, capazes de expressar sensações, sentimentos, pensamentos e de desenvolver seus próprios percursos criativos, articulando a percepção, a sensibilidade, a imaginação, a cognição, sob a orientação do professor.
O trabalho nesse Campo de Experiências deve permitir a imersão das \textbf{crianças} nas diferentes linguagens artísticas e a familiarização e apropriação de várias formas de expressão ( gestual, verbal, \textbf{plástica}, \textbf{dramática} e musical ), através de ricas e variadas experiências de conhecimento, apreciação, expressão e interação. [ ^7 ] : Conviver com diferentes manifestações artísticas, culturais e científicas, locais e universais, no cotidiano da instituição escolar, possibilita às \textbf{crianças}, por meio de experiências diversificadas, vivenciar diversas formas de expressão e linguagens, como as artes visuais ( \textbf{pintura}, modelagem, colagem, fotografia etc. ), a música, o teatro, a dança e o audiovisual, entre outras.
Com base nessas experiências, elas se expressam por várias linguagens, criando suas próprias produções artísticas ou culturais, exercitando a autoria ( coletiva e individual ) com \textbf{sons}, traços, \textbf{gestos}, danças, \textbf{mímicas}, encenações, canções, desenhos, modelagens, manipulação de diversos materiais e de recursos tecnológicos.
Essas experiências contribuem para que, desde muito \textbf{pequenas}, as \textbf{crianças} desenvolvam senso estético e crítico, o conhecimento de si mesmas, dos outros e da realidade que as cerca.
Portanto, a Educação \textbf{Infantil} precisa promover a participação das \textbf{crianças} em tempos e espaços para a produção, manifestação e apreciação artística, de modo a favorecer o desenvolvimento da sensibilidade, da criatividade e da expressão pessoal das \textbf{crianças}, permitindo que se apropriem e reconfigurem, permanentemente, a cultura e potencializem suas singularidades, ao ampliar repertórios e interpretar suas experiências e vivências artísticas. ( BNCC, 2017, p. 39 ).

% matched lemmas: acompanhamento, areia, cheio, criança, dramático, gesto, infantil, mímico, pequeno, pintura, plástico, som, textura, tinta
\end{textsample}
